\section{Introduction}
\label{sec:introduction}

% Draft for Ammaar's revision. Voice: measured, claims scoped to what
% the artifacts support. Every number is a generated macro.

\IEEEPARstart{H}{and-drawn} circuit diagrams remain the lingua franca
of electronics education, whiteboard design discussions, and lab
notebooks. Converting a photograph of such a sketch into a
machine-readable SPICE netlist would connect this informal medium
directly to simulation, checking, and documentation tools. The task is
deceptively hard: it requires detecting freehand component symbols,
tracing noisy, gapped, and crossing pencil strokes into electrical
nets, associating component terminals with those nets, and finally
emitting a netlist that a simulator will actually accept.

Recent systems report strong headline results on this task
\cite{wang2026topology,sina2026,reddy2022handdrawn}, yet the field has
a reproducibility problem: to our knowledge, \emph{no public
evaluation benchmark with human-verified connectivity ground truth
exists}. Wang et al.\ \cite{wang2026topology} train on the public
Digitize-HCD dataset \cite{digitizehcd_data} but evaluate on an
independent 1{,}317-image benchmark that is not released; SINA
\cite{sina2026} and Netlistify \cite{netlistify} evaluate on their own
curated sets. Published numbers are therefore not comparable, and none
of the underlying ground truth is available for scrutiny. Separately,
most pipelines stop at topology: whether the emitted netlist actually
\emph{simulates} — and what silent assumptions were needed to make it
simulate — is rarely measured or disclosed.

This paper addresses both gaps. Our contributions are:

\begin{enumerate}
  \item \textbf{C1 — the first published evaluation benchmark on
  Digitize-HCD.} Frozen train/val/test splits (895/192/\NumTestImages),
  human-verified ground-truth topology for every test image, a strict
  GT validator, and released code that regenerates every number in
  this paper. Digitize-HCD has been used for \emph{training}
  \cite{wang2026topology}; nobody has published an evaluation protocol
  or baseline results on it.

  \item \textbf{C2 — structural connectivity recovery that measurably
  beats classical CV.} A progression of interpretable connectivity
  stages — binarized-ink wire extraction, stitching of mask-induced
  wire gaps, and crossover-aware net assembly driven by the detector's
  \emph{Wire Crossover} class — lifts net-level $F_1$ from
  \AblNetFOneClassical{} to \AblNetFOneVFive{} and terminal-pair $F_1$
  from \AblTpFOneClassical{} to \AblTpFOneVFive{}, with every step
  validated by a paired per-image bootstrap test.

  \item \textbf{C3 — port-aware terminal handling.} Class-aware
  terminal semantics (pin identity and polarity for directional
  two-terminal parts) derived from Digitize-HCD's port-location
  annotations, an annotation modality no published work exploits.
  \draftnote{scope this bullet to exactly what lands in
  Sec.~\ref{sec:method}; MSP = two-terminal polarity + pin order.}

  \item \textbf{C4 — metric standardization and stage attribution.} A
  metric cascade from detection mAP through terminal-pair/net $F_1$,
  normalized graph edit distance, strict end-to-end success, SPICE
  validity, and DC solvability, reported together with bootstrap
  confidence intervals; plus a ground-truth-injection oracle that
  attributes end-to-end error to individual pipeline stages.

  \item \textbf{C5 — transparent, minimal design-intent repair.} An
  electrical-rule-check layer diagnoses why a recovered circuit cannot
  be DC-solved, and a repair layer applies the \emph{minimal} set of
  explicit fixes, each logged in a human-reviewable assumption ledger
  and classified as a behavior-invariant \emph{gauge} choice or a
  flagged, behavior-changing \emph{assumption} with alternatives.
  Repair raises DC solvability from \RepSolvBefore{} to
  \RepSolvAfter{} while provably never altering the recovered
  topology.
\end{enumerate}

Together these turn an anecdotally-evaluated task into a measured one:
the benchmark quantifies exactly where current pipelines fail (wire
topology and terminal association, not symbol detection — detection is
essentially solved at \DetMapFifty{} mAP@0.5), and the repair layer
makes the final, usually-hidden step from ``topologically plausible''
to ``simulable'' explicit and auditable.

\draftnote{closing paragraph: one-sentence paper roadmap.}
